\section{Reduction to a smooth link and its suspension}\label{red:section}\label{sec:reductions}

We first record the consequences of energy minimality needed to put the
argument on a smooth spherical domain.  Throughout Section~\ref{sec:reductions}, a minimizer
has the full ballwise Sobolev comparison property defined in
Section~\ref{sec:introduction}.

\subsection{The regularity inputs}\label{red:regularity-inputs}

The ballwise convention also gives comparison on every bounded open set
$O\Subset\R^d$: extend an $H^1_0(O)$ competitor difference by zero into a ball
containing $\overline O$, apply minimality there, and cancel the energy outside
$O$.  Conversely, comparison on bounded open sets includes comparison on
balls.  It is even enough to have comparison against \emph{arbitrary}
admissible competitors in every origin-centered ball, by the same extension
argument.  There is no restriction to homogeneous competitors.

The convention that competitors are nonnegative makes no difference to the
standard one-phase results.  If $v-u\in H^1_0(O)$ and $u\geq0$, then
$v_+-u\in H^1_0(O)$ and
\[
 \int_O\bigl(|\nabla v_+|^2+\mathbf 1_{\{v_+>0\}}\bigr)
 \leq
 \int_O\bigl(|\nabla v|^2+\mathbf 1_{\{v>0\}}\bigr).
\]
For completeness, the assertion about $H^1_0(O)$ follows by approximating
$v-u$ in $H^1$ by compactly supported smooth functions and applying the
continuity of positive-part truncation in $H^1$; the truncated differences
in the approximation have compact support in $O$.

\begin{proposition}[Regularity and blowups of minimizers]\label{red:regularity}
Let $d\geq2$ and let $u\geq0$ be a local minimizer of the one-phase energy
with volume coefficient $1$ in an open set $D\subset\R^d$.
\begin{enumerate}
\renewcommand{\labelenumi}{\textup{(\roman{enumi})}}
\item There is a locally Lipschitz representative of $u$, and it is harmonic
in $\{u>0\}$.
\item If $x_0\in D\cap\partial\{u>0\}$, every sequence $r_j\downarrow0$
has a subsequence for which
\begin{equation}\label{red:blowup-sequence}
 u_{x_0,r_j}(y):=\frac{u(x_0+r_jy)}{r_j}\longrightarrow U(y)
\end{equation}
locally uniformly in $\R^d$.  The limit is a nonzero, nonnegative,
one-homogeneous global minimizer.  Along a further subsequence the convergence
is strong in $H^1_{\mathrm{loc}}$, and the positivity indicators converge in
$L^1_{\mathrm{loc}}$.
\item If one such limit equals $(y\cdot e)_+$ with $|e|=1$, then near
$x_0$ the free boundary is a smooth hypersurface and $\{u>0\}$ is exactly
one of its sides.  The function $u$ is smooth up to that hypersurface from
the positive side, and
\begin{equation}\label{red:bernoulli}
 u=0,\qquad \nabla u=N_E
 \quad\hbox{on the free boundary},
\end{equation}
where $N_E$ is the Euclidean unit normal pointing into $\{u>0\}$.
\end{enumerate}
\end{proposition}

\begin{proof}[Sources and applicability]
We use the constant-coefficient theory with $\Lambda=1$ in
\cite{Velichkov2023}.  For each ball $B\Subset D$, the restriction of $u$
to $B$ minimizes against every $H^1_0(B)$ competitor difference.
Zero extension of differences therefore gives comparison on every
$O\Subset B$, as required by Definitions~2.11--2.12 there. We apply the
local results inside such balls. For global minimizers, the containing-ball
argument above gives the comparison convention on all bounded open sets.
Local Lipschitz continuity is Theorem~3.1. Harmonicity also follows
directly: for a smooth test supported compactly in the positive phase, both
signs of a sufficiently small additive variation preserve that phase, so the
Dirichlet first variation vanishes.

The Lipschitz bound and $u(x_0)=0$ give locally uniform compactness of the
rescalings.  Proposition~4.1 gives
$\sup_{B_r(x_0)}u\geq\kappa r$ for small $r$, with $\kappa>0$, so a limit
cannot vanish identically.  Proposition~6.2 gives strong Sobolev and
positivity-indicator convergence and the full minimizing property of the
limit; Proposition~9.12 gives its one-homogeneity.  All these statements can
be applied in a bounded ball compactly contained in $D$ before taking the
rescalings.  In particular, the assertion about minimality is the minimizing
compactness result, including the volume term.

Definition~6.10 and Corollary~8.2, with the one-sided graph conclusion in
Theorem~8.1, give a $C^{1,\alpha}$ regular neighborhood from a single
half-space blowup.  Proposition~7.1 supplies the viscosity Bernoulli condition.

To justify all subsequent boundary derivatives, use
\cite[Theorem~1.1 and Corollary~1.2]{DeSilvaFerrariSalsa2019} with
$f_+=f_-=0$ and $G(b)=1+b$.  We check the two-phase viscosity condition
of \cite[Definition~7.1]{DeSilvaFerrariSalsa2014} when $u\geq0$.
Let $v$ be a comparison touching $u$ at a free-boundary
point $x_0$, with a $C^2$ interface and $C^2$ branches having nonnegative
normal slopes $a$ and $b$ on its positive and negative sides.  If $v$ touches
from below, then $v_+\leq u$ in a neighborhood.  For $a>0$, extend its
positive branch to a signed $C^2$ function $\phi$, negative just across
the interface; then $\phi\leq u$ and touches at $x_0$.  The one-phase
condition gives $a\leq1\leq G(b)$; the case $a=0$ is immediate.  If $v$
touches from above, the neighborhood inequality $v\geq u\geq0$ forces
$v=u=0$ throughout its nonpositive side, so $b=0$.  A $C^2$ extension
$\phi$ of its positive branch satisfies $\phi_+\geq u$ on both sides
and touches at $x_0$.  Thus the one-phase condition gives $a\geq1=G(0)$.
These uses of the one-phase condition are justified by
\cite[Definition~7.6 and Proposition~7.1]{Velichkov2023}: to meet its
smooth-test convention, replace each $C^2$ extension by its quadratic
Taylor polynomial minus, respectively plus, $\varepsilon|x-x_0|^2$.
For any $\varepsilon>0$, these are smooth lower and upper supports in a
sufficiently small neighborhood, with the same contact value and gradient.

The interior equations hold because $u$ is harmonic where positive and
identically zero on the interior of its zero set.  The function $G$ is smooth,
strictly increasing and unbounded, with $G(0)=1$; the cited results allow
the negative phase to vanish.  The preceding $C^{1,\alpha}$ one-sided graph
satisfies the slab flatness hypothesis after rotation and sufficiently small
dilation, while the local Lipschitz bound is preserved.  Theorem~1.1 gives
$C^{2,\gamma}$ regularity, and Corollary~1.2 gives smoothness.  Local harmonic boundary regularity with zero Dirichlet data
\cite[Theorems~4.30--4.31 and Corollary~4.32]{Hunter2014} then gives
smoothness of $u$ from the positive side. These estimates apply after
flattening a smooth boundary patch and localizing away from its artificial
boundary; iteration and Sobolev embedding give every required derivative.  The Bernoulli condition and nonnegativity give
\eqref{red:bernoulli}.
\end{proof}

As in Section~\ref{sec:introduction}, a free-boundary point is regular
if it has a half-space blowup and singular otherwise.  Proposition~\ref{red:regularity} explains the
geometric meaning of regularity.  Conversely, a smooth one-sided free
boundary with \eqref{red:bernoulli} has a half-space blowup by its first-order
Taylor expansion.  We henceforth use the continuous representative of a
minimizer.  If it is one-homogeneous almost everywhere, it is one-homogeneous
everywhere: for each fixed $r>0$, both sides of $u(rx)=ru(x)$ are continuous
and agree almost everywhere.

\subsection{Dimension reduction at a nonzero point}\label{red:dimension-reduction}

The next two lemmas give the dimension-reduction steps explicitly;
compare \cite[Lemmas~10.9--10.10]{Velichkov2023}.

\begin{lemma}[Translation invariance of an off-origin blowup]\label{red:translation}
Let $u$ be a one-homogeneous global minimizer, let
$0\ne x_0\in\partial\{u>0\}$, and let $U$ be any limit in
\eqref{red:blowup-sequence}.  Then
\[
 U(y+t x_0)=U(y)\qquad(y\in\R^d,\ t\in\R).
\]
\end{lemma}

\begin{proof}
Write $u_r=u_{x_0,r}$.  For fixed $t\in\R$ and small $r>0$, set
$r'=r/(1+rt)>0$.  Homogeneity gives the exact identity
\begin{equation}\label{red:rescaling-identity}
 u_r(y+t x_0)
 =\frac{u\bigl((1+rt)(x_0+r'y)\bigr)}r
 =u_{r'}(y).
\end{equation}
Let $L$ be a Lipschitz constant for $u$ on a fixed neighborhood of $x_0$.
Since $u(x_0)=0$, for $|y|\leq M$ and small $r$,
\begin{align*}
 |u_{r'}(y)-u_r(y)|
 &\leq\frac{L|r-r'|M}{r'}
       +LrM\left|\frac1{r'}-\frac1r\right|\\
 &=2LM\frac{|r-r'|}{r'}=2LM|t|r.
\end{align*}
Thus the $r'$ rescalings along the same chosen sequence have the same
locally uniform limit as the $r$ rescalings.  Passing to that limit in
\eqref{red:rescaling-identity} proves the assertion for every fixed $y,t$.
No uniqueness of blowups is required.
\end{proof}

\begin{lemma}[Minimality of a translation-invariant factor]\label{red:factor-minimality}
Suppose $U(y,s)=v(y)$ is a global minimizer in
$\R^k\times\R$, where $v\geq0$ belongs to $H^1_{\mathrm{loc}}(\R^k)$.
Then $v$ is a global minimizer in $\R^k$.
\end{lemma}

\begin{proof}
Suppose that a ball $B\subset\R^k$ admits a nonnegative
$q\in H^1(B)$ with $q-v\in H^1_0(B)$ and
\[
 J(v;B)-J(q;B)=\delta>0.
\]
Let $\psi$ be the zero extension of $q-v$.  For $L>1$, choose
$\eta_L:\R\to[0,1]$ equal to $1$ on $[-L,L]$, equal to $0$ outside
$[-L-1,L+1]$, and linear on each intervening interval.  Define
\[
 V_L(y,s)=U(y,s)+\eta_L(s)\psi(y).
\]
This function is nonnegative: on $B\times\R$ it is the convex combination
$(1-\eta_L)v+\eta_Lq$.  Its difference from $U$ belongs to $H^1$ and has
compact support in a sufficiently large ball $\mathcal B_L\subset\R^{k+1}$,
so it belongs to $H^1_0(\mathcal B_L)$.

The energy difference on $B\times[-L,L]$ is $-2L\delta$.  On either
transition cylinder, convexity gives
\[
 |\nabla_y V_L|^2\leq(1-\eta_L)|\nabla v|^2+\eta_L|\nabla q|^2,
 \qquad |\partial_sV_L|^2\leq|q-v|^2.
\]
The difference of the positivity indicators is at most $1$.  Hence the
sum of the two transition contributions is bounded above by
\[
 C=2\int_B\bigl(|\nabla q|^2+|q-v|^2+1\bigr),
\]
independently of $L$.  Outside these three cylinders the energy is unchanged.
Thus
\[
 J(V_L;\mathcal B_L)-J(U;\mathcal B_L)\leq-2L\delta+C<0
\]
for large $L$, contradicting minimality of $U$.
\end{proof}

In dimension one, every nonnegative one-homogeneous Sobolev function
has the form $u(x)=b_+x_++b_-(-x)_+$ with $b_+,b_-\ge0$.
If both slopes are positive, a nonnegative smooth bump $\phi$ supported
in $(-R,R)$ with $\phi(0)>0$ preserves positivity volume and gives
\[
 J(u+\tau\phi;(-R,R))-J(u;(-R,R))
 =-2\tau(b_++b_-)\phi(0)+\tau^2\int_{-R}^R|\phi'|^2<0
\]
for small $\tau>0$. Thus a nonzero minimizing cone has just one positive
slope. After reflection write $u(x)=b x_+$, $b>0$.
The competitors $v_t(x)=bR(x-t)_+/(R-t)$ on $(-R,R)$, $|t|<R$,
have the same trace and energy $b^2R^2/(R-t)+(R-t)$.
Its derivative at the minimizing parameter $t=0$ is $b^2-1=0$,
so $b=1$. This proves flatness in dimension one directly.

The lower-dimensional classification of Jerison--Savin
\cite[Theorem~1.1 and Corollary~1.2]{JerisonSavin2015} and the
dimension-seven minimizing example of De Silva--Jerison
\cite[Theorem~1.1]{DeSilvaJerison2009} give $5\leq d_*\leq7$.
De Silva--Jerison's Theorem~1.1 uses the same functional and arbitrary
ballwise Dirichlet competitors.  For the exclusion argument, suppose from
now on that
\begin{equation}\label{red:first-dimension}
 d=d_*\in\{5,6\},
\end{equation}
and fix a nonzero nonflat one-homogeneous global minimizer $u$ in $\R^d$.

\begin{proposition}[Smooth original link]\label{red:original-link}
Every point of $\partial\{u>0\}\setminus\{0\}$ is regular.  Put
$m_o=d-1$ and $g_o=u|_{S^{m_o}}$.  The set $\{g_o>0\}$ is nonempty and
proper, has finitely many connected components, and their closures are
pairwise disjoint compact smooth manifolds with boundary.  On each such
component $\Omega_o$, the restriction of $g_o$ is smooth up to the boundary
from within $\Omega_o$ and satisfies
\begin{equation}\label{red:link-equations}
 \Delta_{S^{m_o}}g_o=-m_o g_o\quad\hbox{in }\Omega_o,
 \qquad g_o=0,\quad \nabla_{S^{m_o}}g_o=N_o
       \quad\hbox{on }\partial\Omega_o,
\end{equation}
where $N_o$ is the inward spherical unit normal.  In particular,
$\partial\Omega_o\ne\varnothing$.
\end{proposition}

\begin{proof}
At a nonzero free-boundary point $x_0$, choose a blowup $U$ from
Proposition~\ref{red:regularity}.  By Lemma~\ref{red:translation}, in
orthogonal coordinates $\R^d=x_0^\perp\times\R$ it has the form
$U(y,s)=v(y)$.  The factor $v$ is locally Lipschitz, nonnegative, nonzero
and one-homogeneous.  Lemma~\ref{red:factor-minimality} shows that it is a
global minimizer.  Since its dimension is $d-1<d_*$, it is flat.  Therefore
$U$ itself is a half-space solution, and Proposition~\ref{red:regularity}
makes $x_0$ regular.

The spherical positive set is nonempty by nontriviality.  The polar-coordinate
formula for the Euclidean Laplacian of $u(ry)=r g_o(y)$ gives
\begin{equation}\label{red:polar-laplacian}
 \Delta_{\R^d}u(ry)=r^{-1}
       \bigl(\Delta_{S^{m_o}}g_o(y)+m_o g_o(y)\bigr)
\end{equation}
in the positive phase.  If $g_o>0$ on the entire sphere, integrating the
harmonicity identity over the sphere would give $m_o\int g_o=0$, a
contradiction.  Thus its positive set is proper.

The conical free boundary has the radial direction tangent to it at every
nonzero point: the curve $r\mapsto rx$ lies in that free boundary.
Consequently its smooth hypersurface is transverse to the unit sphere,
and their intersection is a smooth hypersurface of $S^{m_o}$.  Its regular
charts have the spherical positive phase exactly on one side.  Each point of
the sphere therefore has a neighborhood whose intersection with $\{g_o>0\}$
is either connected or empty (using a small ball for points away from the
free boundary).  A finite subcover of the sphere shows that there are only
finitely many positive components.  If two component closures met, their
common point would be a free-boundary point; the connected positive side of
a regular chart there would meet both components, forcing them to coincide.
The same charts show that each closure is a smooth compact domain.  Each
component has nonempty boundary because the sphere is connected and the
positive set is proper.

Smoothness of $u$ from its positive side gives smoothness of $g_o$ on these
closures.  Equation~\eqref{red:polar-laplacian} gives the interior equation
in \eqref{red:link-equations}.  On the boundary, the radial derivative of
$u$ is $g_o=0$, so the Euclidean gradient is tangent to the sphere.
Equation~\eqref{red:bernoulli} then gives the remaining boundary conditions.
\end{proof}

\begin{lemma}[Selection of a nonflat component]\label{red:nonflat-component}
There is a component $\Omega_o$ in Proposition~\ref{red:original-link}
such that
\begin{equation}\label{red:original-hessian}
 T_o:=\nabla^2_{S^{m_o}}g_o+g_o\Id\not\equiv0
 \quad\hbox{on }\Omega_o.
\end{equation}
Here $\Id$ denotes the spherical metric, viewed as a symmetric two-tensor.
\end{lemma}

\begin{proof}
For $y$ in the spherical positive set and $r>0$,
\[
 \nabla_{\R^d}u(ry)=g_o(y)y+\nabla_{S^{m_o}}g_o(y).
\]
Differentiating this formula, or using Euler's identity for the radial
entries, gives
\begin{equation}\label{red:ambient-hessian}
 D^2u(ry)[y,\cdot]=0,
 \qquad
 D^2u(ry)[v_1,v_2]=r^{-1}T_o(y)[v_1,v_2]
 \quad(v_1,v_2\perp y).
\end{equation}
Suppose that $T_o$ vanished identically on every positive component.  Then
$D^2u=0$ on each corresponding connected open cone $D_i$, so
$u(x)=b_i\cdot x+c_i$ there.  Homogeneity forces $c_i=0$, and positivity
forces $b_i\ne0$.  Thus $D_i\subset H_i:=\{b_i\cdot x>0\}$.
The set $D_i$ is also closed relative to $H_i$: if
$x\in\overline{D_i}\cap H_i$, continuity gives $u(x)=b_i\cdot x>0$;
a small positive neighborhood of $x$ meets $D_i$, hence belongs to the
same component.  Since $H_i$ is connected, $D_i=H_i$.
The Bernoulli condition at any nonzero point of its boundary gives
$|b_i|=1$.

Two disjoint open half-spaces through the origin must have opposite unit
normals: if their normals are $e,f$ with $f\ne-e$, the vector $e+f$ lies
in both. But opposite half-spaces have spherical closures meeting along
the equator $e^\perp\cap S^{m_o}$, contrary to the pairwise disjoint
closures in Proposition~\ref{red:original-link}. Thus there is exactly
one positive component, and $u=(x\cdot e)_+$, contrary to nonflatness.
This proves \eqref{red:original-hessian}.
\end{proof}

\subsection{Suspension and the working domain}\label{red:suspension-subsection}

\begin{proposition}[Suspension preserves global minimality]\label{red:suspension}
If $u\geq0$ is a global minimizer in $\R^k$, then
$\widetilde u(x,s)=u(x)$ is a global minimizer in $\R^{k+1}$.
If $u$ is the cone in \eqref{red:first-dimension} with $d=5$, the singular
free-boundary set of $\widetilde u$ is exactly $\{0\}\times\R$.
For the component chosen in Lemma~\ref{red:nonflat-component}, the
corresponding positive component on $S^5$ is
\begin{equation}\label{red:suspension-coordinates}
 \widetilde\Omega
 =\bigl\{\Phi(t,y):=(\sin t\,y,\cos t):
              0<t<\pi,\ y\in\Omega_o\bigr\},
 \qquad
 \widetilde g(\Phi(t,y))=\sin t\,g_o(y).
\end{equation}
Its boundary consists of the smooth side
$\Phi((0,\pi)\times\partial\Omega_o)$ and the two singular tips
$p_+=(0,1)$ and $p_-=(0,-1)$.
\end{proposition}

\begin{proof}
Fix a ball $\mathcal B=B_R((x_c,s_c))\subset\R^{k+1}$ and a nonnegative
$V\in H^1(\mathcal B)$ with $V-\widetilde u\in H^1_0(\mathcal B)$.
Extend that difference by zero to $F\in H^1(\R^{k+1})$.  For
$|s-s_c|<R$, let
\[
 B_s=B_{\sqrt{R^2-(s-s_c)^2}}(x_c)
\]
be the horizontal slice.  For almost every such $s$,
$F(\cdot,s)|_{B_s}\in H^1_0(B_s)$.  To see the zero-boundary assertion
without a trace regularity assumption on $V$, approximate $F$ in
$H^1(\R^{k+1})$ by the zero extensions of functions in
$C_c^\infty(\mathcal B)$.  After passing to a subsequence, Fubini's Theorem
gives convergence of their slices to $F(\cdot,s)$ in $H^1(\R^k)$ for
almost every $s$.  Each approximating slice lies in $C_c^\infty(B_s)$,
so the assertion follows by closure.

For almost every $s$, $V_s=V(\cdot,s)$ is therefore an admissible
nonnegative competitor for $u$ in $B_s$.  Integrating its minimizing
inequality and using the nonnegativity of the additional derivative energy,
we obtain
\begin{align*}
 J(\widetilde u;\mathcal B)
 &=\int_{s_c-R}^{s_c+R}J(u;B_s)\,ds\\
 &\leq\int_{s_c-R}^{s_c+R}J(V_s;B_s)\,ds
 \leq J(V;\mathcal B).
\end{align*}
This proves the full ballwise minimizing property, for arbitrary centers
and arbitrary admissible Sobolev competitors.  One-homogeneity and
nontriviality are plainly preserved.  Flatness of $\widetilde u$ would
force its half-space normal to have zero last coordinate, and would
therefore force $u$ to be flat.

In the conical case, $u(0)=0$ and a positive ray approaches the origin,
so $0\in\partial\{u>0\}$.  The product identity
\[
 \partial\{\widetilde u>0\}=\partial\{u>0\}\times\R
\]
shows that all its off-axis points are regular: their positive phase is
the product of a smooth one-sided regular chart for $u$ and an interval.
At any axis point $z_0=(0,s_0)$, every rescaling is exactly
\[
 \frac{\widetilde u(z_0+r(x,s))}{r}
 =\frac{u(rx)}r=\widetilde u(x,s).
\]
This is nonflat, so no axis point is regular.  The axis meets $S^5$
only in $p_+$ and $p_-$.

Away from these two points, $\Phi$ is a smooth product coordinate map
from $(0,\pi)\times S^4$ to $S^5\setminus\{p_+,p_-\}$.
Positivity in these coordinates is precisely $g_o(y)>0$, so connected
components are as in \eqref{red:suspension-coordinates}.  Since
$\overline{\Omega_o}$ is compact, its suspended closure is obtained by
adding the two endpoints of the interval, with each endpoint collapsed
to its tip.  Its remaining boundary is exactly the stated smooth side.
\end{proof}

\paragraph{Notation for the remaining argument.}\label{red:working-domain}
Fix the original component $\Omega_o\subset S^{m_o}$ selected in
Lemma~\ref{red:nonflat-component}, where $m_o\in\{4,5\}$.  If $m_o=5$,
set $\Omega=\Omega_o$, $g=g_o$, and $\Sigma=\varnothing$.  If $m_o=4$,
set $\Omega=\widetilde\Omega$, $g=\widetilde g$, and
$\Sigma=\{p_+,p_-\}$.  In either case $\Omega\subset S^5$ is connected,
and $g$ comes from a nonnegative one-homogeneous global minimizer in
$\R^6$.  The closure is smooth away from $\Sigma$, with positive phase
on one side of its regular boundary
\[
 \partial_{\mathrm{reg}}\Omega:=\partial\Omega\setminus\Sigma.
\]
All spherical boundary values below are taken from the positive side.
Tests on $\overline\Omega$ will be smooth and supported away from $\Sigma$.
Every compact portion of the selected Euclidean cone's closure that avoids
the origin and, in the suspended case, the axis is separated from all other
positive components.  This follows from the disjoint compact original
closures, and in the suspension from the same fact on any closed subinterval
of $(0,\pi)$.  Consequently the local variations on this component can be
made as variations of the global minimizer.  No minimizing property of the
selected component extended by zero is needed.
